% GENERATED FILE. Edit canonical lessons, then run npm run generate:books.
% canonical-source-hash: fnv1a64:b9fa7ad38bd842ff
% canonical-lessons: TE-C121-urumu, TE-C121-merupu, TE-C121-tupanu, TE-C121-ala, TE-C121-konda

\chapter{Thunder, Lightning, Storm, Wave, Hill}
\label{ch:te-thunder-lightning-storm-wave-hill}
\hlchaptermodality{121}

\begin{chapteropening}
\textbf{By the end of this chapter:} \emph{I can say thunder, lightning, a storm, a wave and a hill.}

Complete the last lesson of chapter 121: I can say thunder, lightning, a storm, a wave and a hill.
\end{chapteropening}

\section[\texorpdfstring{\te{ఉరుము} (urumu)}{ఉరుము (urumu)}]{\te{ఉరుము} (urumu) — thunder}

\label{lesson:TE-C121-urumu}



\begin{quote}
Before the new one: say the Telugu for a leaf, then the Telugu for soil.
\end{quote}

\subsection*{You'll want to know: \te{ఉరుము}}
\textbf{\te{ఉరుము}} — \emph{urumu} — \textquotedblleft{}thunder\textquotedblright{}.

\begin{grammarlens}[title={What you've built}]
One more word for this chapter.
\end{grammarlens}

\subsection*{Your turn}
\begin{itemize}
\raggedright
  \item \emph{Say it:} \emph{urumu}
  \item \emph{Say it:} \emph{urumu}, once more
  \item \emph{Say it:} say \emph{urumu}
  \item \emph{Recall:} say the Telugu for a leaf, then the Telugu for soil, then say \emph{urumu} again
\end{itemize}

\begin{tcolorbox}[breakable,colback=teal!4,colframe=teal!35!black,title={Before you move on}]
What does \te{ఉరుము} mean? (\textquotedblleft{}Thunder\textquotedblright{}.) Say it once more.
\end{tcolorbox}
\section[\texorpdfstring{\te{మెరుపు} (merupu)}{మెరుపు (merupu)}]{\te{మెరుపు} (merupu) — lightning}

\label{lesson:TE-C121-merupu}



\begin{quote}
Before the new one: say the Telugu for soil, then the Telugu for thunder.
\end{quote}

\subsection*{You'll want to know: \te{మెరుపు}}
\textbf{\te{మెరుపు}} — \emph{merupu} — \textquotedblleft{}lightning\textquotedblright{}.

\begin{grammarlens}[title={What you've built}]
One more word for this chapter.
\end{grammarlens}

\subsection*{Your turn}
\begin{itemize}
\raggedright
  \item \emph{Say it:} \emph{merupu}
  \item \emph{Say it:} \emph{merupu}, once more
  \item \emph{Say it:} say \emph{merupu}
  \item \emph{Recall:} say the Telugu for soil, then the Telugu for thunder, then say \emph{merupu} again
\end{itemize}

\begin{tcolorbox}[breakable,colback=teal!4,colframe=teal!35!black,title={Before you move on}]
What does \te{మెరుపు} mean? (\textquotedblleft{}Lightning\textquotedblright{}.) Say it once more.
\end{tcolorbox}
\section[\texorpdfstring{\te{తుపాను} (tupānu)}{తుపాను (tupānu)}]{\te{తుపాను} (tupānu) — a storm}

\label{lesson:TE-C121-tupanu}



\begin{quote}
Before the new one: say the Telugu for thunder, then the Telugu for lightning.

Name the seasons: \textbf{\te{వసంత} \te{ఋతువు}}, \textbf{\te{వేసవి}}, \textbf{\te{వానాకాలం}}, \textbf{\te{శీతాకాలం}}.
\end{quote}

\subsection*{You'll want to know: \te{తుపాను}}
\textbf{\te{తుపాను}} — \emph{tupānu} — \textquotedblleft{}a storm\textquotedblright{}.

\begin{grammarlens}[title={What you've built}]
One more word for this chapter.
\end{grammarlens}

\subsection*{Your turn}
\begin{itemize}
\raggedright
  \item \emph{Say it:} \emph{tupānu}
  \item \emph{Say it:} \emph{tupānu}, once more
  \item \emph{Say it:} say \emph{tupānu}
  \item \emph{Recall:} say the Telugu for thunder, then the Telugu for lightning, then say \emph{tupānu} again
\end{itemize}

\begin{tcolorbox}[breakable,colback=teal!4,colframe=teal!35!black,title={Before you move on}]
What does \te{తుపాను} mean? (\textquotedblleft{}A storm\textquotedblright{}.) Say it once more.
\end{tcolorbox}
\section[\texorpdfstring{\te{అల} (ala)}{అల (ala)}]{\te{అల} (ala) — a wave}

\label{lesson:TE-C121-ala}



\begin{quote}
Before the new one: say the Telugu for lightning, then the Telugu for a storm.
\end{quote}

\subsection*{You'll want to know: \te{అల}}
\textbf{\te{అల}} — \emph{ala} — \textquotedblleft{}a wave\textquotedblright{}.

\begin{grammarlens}[title={What you've built}]
One more word for this chapter.
\end{grammarlens}

\subsection*{Your turn}
\begin{itemize}
\raggedright
  \item \emph{Say it:} \emph{ala}
  \item \emph{Say it:} \emph{ala}, once more
  \item \emph{Say it:} say \emph{ala}
  \item \emph{Recall:} say the Telugu for lightning, then the Telugu for a storm, then say \emph{ala} again
\end{itemize}

\begin{tcolorbox}[breakable,colback=teal!4,colframe=teal!35!black,title={Before you move on}]
What does \te{అల} mean? (\textquotedblleft{}A wave\textquotedblright{}.) Say it once more.
\end{tcolorbox}
\section[\texorpdfstring{\te{కొండ} (koṇḍa)}{కొండ (koṇḍa)}]{\te{కొండ} (koṇḍa) — a hill, a mountain}

\label{lesson:TE-C121-konda}



\begin{quote}
Before the new one: say the Telugu for a storm, then the Telugu for a wave.
\end{quote}

\subsection*{You'll want to know: \te{కొండ}}
\textbf{\te{కొండ}} — \emph{koṇḍa} — \textquotedblleft{}a hill, a mountain\textquotedblright{}.

\begin{grammarlens}[title={What you've built}]
One more word for this chapter.
\end{grammarlens}

\subsection*{Your turn}
\begin{itemize}
\raggedright
  \item \emph{Say it:} \emph{koṇḍa}
  \item \emph{Say it:} \emph{koṇḍa}, once more
  \item \emph{Say it:} say \emph{koṇḍa}
  \item \emph{Recall:} say the Telugu for a storm, then the Telugu for a wave, then say \emph{koṇḍa} again
\end{itemize}

\begin{tcolorbox}[breakable,colback=teal!4,colframe=teal!35!black,title={Before you move on}]
What does \te{కొండ} mean? (\textquotedblleft{}A hill, a mountain\textquotedblright{}.) Say it once more.
\end{tcolorbox}
